

\section{质心：为什么系统还能用一个点代表整体平移}

一个由质量分别为 $m_1, m_2, \cdots, m_n$ 的质点组成的复杂系统所受的合外力 $\vec{F}_{\text{外}}$ 等于总动量 $\vec{P}$ 随时间的变化率：
\[
\vec{F}_{\text{外}} = \frac{\mathrm{d}\vec{P}}{\mathrm{d}t} = \frac{\mathrm{d}}{\mathrm{d}t}\left(\sum_i m_i \vec{v}_i\right)= \frac{\mathrm{d}^2}{\mathrm{d}t^2}\left(\sum_i m_i \vec{r}_i\right)= m \frac{\mathrm{d}^2}{\mathrm{d}t^2}\left(\frac{\sum_i m_i \vec{r}_i}{m}\right)
\]
观察上式，括号内部的矢量代数式 $\frac{\sum_i m_i \vec{r}_i}{m}$ 仅由各个质点的质量以及它们的空间几何位置决定。为此，我们定义质心，其位置矢量 $\vec{r}_C$ ：
\[
\boxed{\vec{r}_C = \frac{\sum_i m_i \vec{r}_i}{m} = \frac{m_1\vec{r}_1 + m_2\vec{r}_2 + \cdots + m_n\vec{r}_n}{m_1 + m_2 + \cdots + m_n}}
\]
若改写成直角坐标分量，就是
 \(x_C=\frac{\sum_i m_i x_i}{M},\qquad y_C=\frac{\sum_i m_i y_i}{M},\qquad z_C=\frac{\sum_i m_i z_i}{M}.\) 

\begin{exercise}[任意三角形三个等质量点的质心]
任意三角形的三个顶点分别放有质量都为 \(m\) 的质点，三点坐标依次为 \((x_1,y_1)\)、\((x_2,y_2)\)、\((x_3,y_3)\)。求系统质心坐标。
\end{exercise}
\begin{solution}
\[
\begin{aligned}
x_C&=\frac{mx_1+mx_2+mx_3}{3m}
=\frac{x_1+x_2+x_3}{3},~y_C=\frac{my_1+my_2+my_3}{3m}
=\frac{y_1+y_2+y_3}{3}.
\end{aligned}
\]
\end{solution}

当研究对象不再是离散的质点系，而是质量连续分布的实体，我们在物体内部任意截取一个微元，将其质量记为 $\mathrm{d}m$，该微元在直角坐标系中的位置矢量为 $\vec{r}$。整个连续体的总质量为 $m = \int \mathrm{d}m$。此时，连续体的质心位置矢量和分坐标形式定义为：$$\vec{r}_C = \frac{\int \vec{r} \mathrm{d}m}{m}~~~x_C = \frac{\int x \mathrm{d}m}{m}~~~y_C = \frac{\int y \mathrm{d}m}{m}~~~ z_C = \frac{\int z \mathrm{d}m}{m}$$

\begin{example}
一段质量为 $m$、半径为 $R$ 的均匀铁丝被弯成半圆形。求此半圆形铁丝的质心位置。
\end{example}
\begin{solution}
以半圆的圆心为坐标原点 $O$，半圆底边所在的直线为 $x$ 轴，对称轴为 $y$ 轴建立平面直角坐标系。由于该均匀半圆铁丝关于 $y$ 轴左右严格对称，根据几何对称性定理，其质心必然位于 $y$ 轴上，故其 $x$ 方向的坐标分量可以直接得出：$x_C = 0$，只需算 $y_C$ 的积分值。
在半圆铁丝上任取一段微元，其对应的圆心角为 $\mathrm{d}\theta$，辐角为 $\theta$。该微元的几何弧长为：$\mathrm{d}l = R\mathrm{d}\theta$，该微元所在的位置分量为：$y = R\sin\theta$，由于铁丝质量分布均匀，我们引入线密度 $\lambda$。半圆铁丝的总长度为 $\pi R$，故线密度可表示为 $\lambda = \frac{m}{\pi R}$。该微元的质量 $\mathrm{d}m$ 满足： \(\mathrm{d}m = \lambda \mathrm{d}l = \frac{m}{\pi R} \cdot R\mathrm{d}\theta = \frac{m}{\pi}\mathrm{d}\theta\) 则
$$y_C = \frac{\int y \mathrm{d}m}{m} = \frac{\int_{0}^{\pi} (R\sin\theta) \cdot \left(\frac{m}{\pi}\mathrm{d}\theta\right)}{m}=\frac{R}{\pi} \int_{0}^{\pi} \sin\theta \mathrm{d}\theta = \frac{R}{\pi} \Big[ -\cos\theta \Big]_0^\pi = \frac{2R}{\pi}$$
\end{solution}

\begin{exercise}
有质量为 $2m$ 的弹丸，从地面斜抛出去，在不受干扰情况下的落地点为 $x_C$。如果它在飞行到最高点处突然爆炸成质量相等的两碎片。其中一碎片铅直自由下落，另一碎片水平抛出，它们同时落地。问第二块碎片落在何处？
\end{exercise}
\begin{solution}
质点系的质心运动完全由系统所受的外力决定，而与系统内部质点之间的相互作用（内力）无关。因此无论弹丸在空中分裂成多少碎片，整个系统的质心运动轨迹依然会保持原有的抛物线不变，落地点依旧是未爆炸前的 $x_C$。由于两碎片同时落地，在它们落地的这一瞬间，直接套用质心的一维分量坐标公式： \(x_C = \frac{m_1 x_1 + m_2 x_2}{m_1 + m_2}\) 代入已知量 $m_1 = m_2 = m$ 以及 $x_1 = 0$：$$x_C = \frac{m \cdot 0 + m \cdot x_2}{m + m} = \frac{m x_2}{2m} = \frac{x_2}{2}\implies x_2 = 2x_C$$
\end{solution}
